% !TEX root = ../main.tex
\section{Ejercicio propuesto 1: implementación y validación del montículo de Fibonacci}

\subsection{Arquitectura y estructura de datos}
El montículo de Fibonacci se modeló mediante dos clases principales: \texttt{FibonacciNode} y \texttt{FibonacciHeap}. Cada nodo mantiene referencias bidireccionales circulares (\texttt{left} y \texttt{right}), punteros de jerarquía (\texttt{parent} y \texttt{child}), su grado (\texttt{degree}) y su bandera booleana de marcado (\texttt{mark}).

La clase \texttt{FibonacciHeap} administra el puntero \texttt{min\_node}, el contador de nodos $n$ y contadores de instrumentación (\texttt{links}, \texttt{cuts}, \texttt{cascading\_cuts}, \texttt{consolidations}, \texttt{roots\_examined}).

\begin{lstlisting}[language=Python, caption={Definición formal de la clase FibonacciNode y constructor de FibonacciHeap}, label={lst:fib_clase}]
class FibonacciNode:
    def __init__(self, key, value=None):
        self.key = key
        self.value = value
        self.parent = None
        self.child = None
        self.left = self
        self.right = self
        self.degree = 0
        self.mark = False

class FibonacciHeap:
    def __init__(self):
        self.min_node = None
        self.n = 0
        self.links = 0
        self.cuts = 0
        self.cascading_cuts = 0
        self.consolidations = 0
        self.roots_examined = 0
\end{lstlisting}

\subsection{Operaciones críticas y manejadores de nodo}
\begin{enumerate}
    \item \textbf{Inserción en $O(1)$:} \texttt{insert(key, value=None)} crea un nodo y lo concatena a la lista circular de raíces junto a \texttt{min\_node}. Retorna directamente el objeto nodo como \textbf{manejador}, permitiendo ejecutar posteriormente \texttt{decrease-key} o \texttt{delete} en tiempo constante sin búsquedas lineales.
    \item \textbf{Extracción del mínimo en $O(\log n)$ amortizado:} \texttt{extract\_min()} promueve todos los hijos del nodo mínimo a la lista de raíces, desvincula \texttt{min\_node} y dispara la consolidación (\texttt{\_consolidate}) mediante un arreglo de grados acotado superiormente por $D(n) \le \lfloor \log_\phi(n) \rfloor + 2$.
    \item \textbf{Disminución de clave y cortes en cascada en $O(1)$ amortizado:} Si la nueva clave viola la propiedad de montículo frente a su padre, se ejecuta \texttt{\_cut}, desvinculando el nodo hacia la raíz. Si el padre ya estaba marcado, se propaga \texttt{\_cascading\_cut} recursivamente hacia arriba.
    \item \textbf{Unión en $O(1)$:} Concatena los punteros de las listas circulares de raíces de dos montículos en tiempo estrictamente constante, sin reinserciones.
\end{enumerate}

\subsection{Verificación exhaustiva de invariantes estructurales}
Para garantizar la solidez de la implementación, el método \texttt{validate()} evalúa rigurosamente las 8 invariantes demandadas por la Guía 04 tras cada operación:

\begin{lstlisting}[language=Python, caption={Método de comprobación automática de invariantes}, label={lst:fib_validador}]
def validate(self, after_extract_min=False):
    if self.min_node is None:
        assert self.n == 0, f"Error: min_node es None pero n={self.n}"
        return True

    roots = self._get_circular_list(self.min_node)
    min_key = min(r.key for r in roots)
    # Invariante 5: Minimo correcto
    assert self.min_node.key == min_key, "Invariante 5 violada: min_node incorrecto"

    # Invariante 7: Raices no marcadas y sin padre
    for r in roots:
        assert not r.mark, f"Invariante 7 violada: raiz {r.key} marcada"
        assert r.parent is None, f"Invariante 3 violada: raiz con padre"

    # Invariante 8: Consolidacion (grados unicos tras extract-min)
    if after_extract_min:
        root_degrees = [r.degree for r in roots]
        assert len(root_degrees) == len(set(root_degrees)), "Grados duplicados"

    # Recorrido exhaustivo de invariantes 1, 2, 3, 4 y 6
    visited = set()
    for r in roots:
        self._validate_node_recursive(r, visited)

    assert len(visited) == self.n, f"Invariante 6: nodos alcanzables ({len(visited)}) != n ({self.n})"
    return True
\end{lstlisting}

\subsection{Resultados de la batería de pruebas obligatorias}
Se ejecutó la batería completa de casos de prueba demandados en las Secciones 9.1 y 16 de la Guía 04. La salida de consola certifica el cumplimiento del 100\% de las aserciones:

\begin{lstlisting}[style=consolestyle, caption={Trazas de ejecución de la batería de pruebas obligatorias}]
>>> Ejecutando Bateria de Pruebas Obligatorias de Invariantes...
  [OK] Inserciones variadas y extraccion ordenada validada.
  [OK] Operacion de Union validada.
  [OK] Disminuciones, cortes, cascadas y eliminaciones validados.
>>> Bateria obligatoria superada con exito (100% asserts pasados).
\end{lstlisting}

\subsubsection*{Casos de prueba evaluados y superados}
\begin{itemize}
    \item \textbf{Extracción sobre estructuras extremas:} Extracción sobre montículo vacío (retorna \texttt{None} sin excepciones no controladas) y extracción sobre montículo de 1 solo nodo.
    \item \textbf{Patrones de inserción:} Inserciones crecientes, decrecientes, aleatorias y con claves repetidas o duplicadas.
    \item \textbf{Extracción total:} Verificación de que la secuencia de claves extraídas es estrictamente monótona no decreciente (\texttt{orden\_extraido == sorted(claves)}).
    \item \textbf{Unión de bosques:} Unión de montículo con otro vacío, y unión de montículos no vacíos sin pérdida de nodos ni violación de \texttt{min\_node}.
    \item \textbf{Disminución y cortes:} Disminución al mismo valor (sin efecto), rechazo formal de clave mayor arrojando \texttt{ValueError}, disminución sin corte, disminución con corte simple, y corte en cascada propagado a más de 2 niveles.
    \item \textbf{Eliminación arbitraria:} Eliminación de nodo raíz y eliminación de nodo interno con múltiples subárboles hijos.
\end{itemize}
